% --- Archon physics formalization source begin ---
% archon:physics
% archon:covers PhyXMiniProblems/problem_phyx_mini_0246.lean
% archon:source-report reports/phyx_mini/problem_phyx_mini_0246.source.json

\chapter{Physics problem phyx\_mini\_0246}
\label{ch:PhyXMiniProblems_problem_phyx_mini_0246}

\paragraph{Problem source.}
\#\# Physical scenario

Two loudspeakers in a plane are \$6.0 \textbackslash{}, \textbackslash{}mathrm\{m\}\$ apart and in phase. They emit equal-amplitude sound waves with a wavelength of \$1.0 \textbackslash{}, \textbackslash{}mathrm\{m\}\$. Each speaker alone creates sound with intensity \$I\_0\$. An observer at point \$A\$ is \$10 \textbackslash{}, \textbackslash{}mathrm\{m\}\$ in front of the plane containing the two loudspeakers and centered between them. A second observer at point \$B\$ is \$10 \textbackslash{}, \textbackslash{}mathrm\{m\}\$ directly in front of one of the speakers.

\#\# Question

In terms of \$I\_0\$, what is the the intensity \$I\_B\$ at point \$B\$?

\#\# Answer choices

- A: A: \$0.95I\_0\$

- B: B: \$1.05I\_0\$

- C: C: \$0.85I\_0\$

- D: D: \$0.75I\_0\$

\#\# Figure caption (auxiliary; use the image as primary evidence)

The image depicts a geometric representation of wave interference. It features two speakers emitting sound waves, positioned vertically apart. Key elements include:

- **Speakers**: Two speakers are shown emitting sound waves, depicted with blue wave lines.
- **Distances**: 
  - The vertical distance between the speakers is labeled as 6.0 m.
  - A horizontal distance marked above the speakers shows 10.0 m.

- **Points A and B**: 
  - Point A is below point B.
  - Lines \textbackslash{}( r\_1 \textbackslash{}) and \textbackslash{}( r\_2 \textbackslash{}) connect each speaker to point B, indicating paths of the sound waves.

- **Wavelength (\textbackslash{}(\textbackslash{}lambda\textbackslash{}))**: The wavelength is marked as 1.0 m for both speakers.
- **Phase Difference (\textbackslash{}(\textbackslash{}Delta \textbackslash{}phi\_0\textbackslash{}))**: Labeled as 0 rad, indicating the waves are in phase.

This diagram likely illustrates a basic wave interference setup, possibly demonstrating concepts like constructive interference.

\#\# Dataset metadata

Resonance and Harmonics; Physical Model Grounding Reasoning, Spatial Relation Reasoning

\paragraph{Recorded answer/context.}
A: A: \$0.95I\_0\$

\paragraph{Figure/image path.}
/root/proposal\_for\_physic/science-mango/phyx\_mini\_run/phyx\_data/test\_image/246.png

\paragraph{Formalization target.}
create a compiling Lean file with sorry bodies at `PhyXMiniProblems/problem\_phyx\_mini\_0246.lean`. The Lean declarations must preserve the physical quantities, dimensions or dimensional roles, figure labels, governing-law hypotheses, and final relation expressed by this problem.
Use Mathlib/Physlib names found through LeanExplore where available. If a physics API is missing, introduce faithful local abstractions rather than scalar placeholder aliases.

\begin{theorem}[Physics formalization target]
\label{thm:physics:phyx_mini_0246:target}
\lean{PhyXMiniProblems.ProblemPhyXMini0246.problem_phyx_mini_0246}
\uses{def:physics:phyx-mini-0246:phyxminiproblems-problemphyxmini0246-acousticintensityreadout, def:physics:phyx-mini-0246:phyxminiproblems-problemphyxmini0246-twoloudspeakersetup, def:physics:phyx-mini-0246:phyxminiproblems-problemphyxmini0246-matchessourcedescription, def:physics:phyx-mini-0246:phyxminiproblems-problemphyxmini0246-matchesproblemreadouts, def:physics:phyx-mini-0246:phyxminiproblems-problemphyxmini0246-matchessuppliedfigure, def:physics:phyx-mini-0246:phyxminiproblems-problemphyxmini0246-hasphysicalacousticparameters, def:physics:phyx-mini-0246:phyxminiproblems-problemphyxmini0246-satisfiesstraightlineraygeometry, def:physics:phyx-mini-0246:phyxminiproblems-problemphyxmini0246-satisfiespathdifferencelaw, def:physics:phyx-mini-0246:phyxminiproblems-problemphyxmini0246-satisfiespropagationphaselaw, def:physics:phyx-mini-0246:phyxminiproblems-problemphyxmini0246-satisfiescoherenttwosourceinterferencelaw, def:physics:phyx-mini-0246:phyxminiproblems-problemphyxmini0246-intensityratioatb, def:physics:phyx-mini-0246:phyxminiproblems-problemphyxmini0246-roundedtonearesthundredth, def:physics:phyx-mini-0246:phyxminiproblems-problemphyxmini0246-isuniquematchingdisplayedintensitychoice}
The assigned autoformalize agent should translate this physics problem into Lean declarations in the covered file, with theorem and lemma proof bodies written as `by sorry`.
\end{theorem}
\begin{proof}
This is an autoformalization task, not a proof task. Produce statements that can later be proved by the physics prover without weakening the physical model.
\end{proof}

% --- Archon declaration topology begin ---
\section{Declaration topology}
\begin{definition}[Length Quantity]
  \label{def:physics:phyx-mini-0246:phyxminiproblems-problemphyxmini0246-lengthquantity}
  \lean{PhyXMiniProblems.ProblemPhyXMini0246.LengthQuantity}
  A nonnegative physical length, independent of its readout unit
\end{definition}
\begin{proof}
  This is the declared model definition of Length Quantity.
\end{proof}
\begin{definition}[Signed Length Quantity]
  \label{def:physics:phyx-mini-0246:phyxminiproblems-problemphyxmini0246-signedlengthquantity}
  \lean{PhyXMiniProblems.ProblemPhyXMini0246.SignedLengthQuantity}
  A signed physical length used for planar coordinates
\end{definition}
\begin{proof}
  This is the declared model definition of Signed Length Quantity.
\end{proof}
\begin{definition}[Acoustic Intensity Quantity]
  \label{def:physics:phyx-mini-0246:phyxminiproblems-problemphyxmini0246-acousticintensityquantity}
  \lean{PhyXMiniProblems.ProblemPhyXMini0246.AcousticIntensityQuantity}
  Acoustic intensity has SI unit watt per square metre and physical dimension mass / time\textasciicircum{}3
\end{definition}
\begin{proof}
  This is the declared model definition of Acoustic Intensity Quantity.
\end{proof}
\begin{definition}[length Readout]
  \label{def:physics:phyx-mini-0246:phyxminiproblems-problemphyxmini0246-lengthreadout}
  \lean{PhyXMiniProblems.ProblemPhyXMini0246.lengthReadout}
  \uses{def:physics:phyx-mini-0246:phyxminiproblems-problemphyxmini0246-lengthquantity}
  Read a nonnegative physical length in a chosen length unit
\end{definition}
\begin{proof}
  This is the declared model definition of length Readout.
\end{proof}
\begin{definition}[signed Length Readout]
  \label{def:physics:phyx-mini-0246:phyxminiproblems-problemphyxmini0246-signedlengthreadout}
  \lean{PhyXMiniProblems.ProblemPhyXMini0246.signedLengthReadout}
  \uses{def:physics:phyx-mini-0246:phyxminiproblems-problemphyxmini0246-signedlengthquantity}
  Read a signed physical coordinate in a chosen length unit
\end{definition}
\begin{proof}
  This is the declared model definition of signed Length Readout.
\end{proof}
\begin{definition}[acoustic Intensity Readout]
  \label{def:physics:phyx-mini-0246:phyxminiproblems-problemphyxmini0246-acousticintensityreadout}
  \lean{PhyXMiniProblems.ProblemPhyXMini0246.acousticIntensityReadout}
  \uses{def:physics:phyx-mini-0246:phyxminiproblems-problemphyxmini0246-acousticintensityquantity}
  Read an acoustic intensity in a coherent system of units
\end{definition}
\begin{proof}
  This is the declared model definition of acoustic Intensity Readout.
\end{proof}
\begin{definition}[length In Meters]
  \label{def:physics:phyx-mini-0246:phyxminiproblems-problemphyxmini0246-lengthinmeters}
  \lean{PhyXMiniProblems.ProblemPhyXMini0246.lengthInMeters}
  \uses{def:physics:phyx-mini-0246:phyxminiproblems-problemphyxmini0246-lengthquantity, def:physics:phyx-mini-0246:phyxminiproblems-problemphyxmini0246-lengthreadout}
  Metre readout of a physical length
\end{definition}
\begin{proof}
  This is the declared model definition of length In Meters.
\end{proof}
\begin{definition}[signed Length In Meters]
  \label{def:physics:phyx-mini-0246:phyxminiproblems-problemphyxmini0246-signedlengthinmeters}
  \lean{PhyXMiniProblems.ProblemPhyXMini0246.signedLengthInMeters}
  \uses{def:physics:phyx-mini-0246:phyxminiproblems-problemphyxmini0246-signedlengthquantity, def:physics:phyx-mini-0246:phyxminiproblems-problemphyxmini0246-signedlengthreadout}
  Metre readout of a signed planar coordinate
\end{definition}
\begin{proof}
  This is the declared model definition of signed Length In Meters.
\end{proof}
\begin{definition}[acoustic Intensity In Watts Per Square Meter]
  \label{def:physics:phyx-mini-0246:phyxminiproblems-problemphyxmini0246-acousticintensityinwattspersquaremeter}
  \lean{PhyXMiniProblems.ProblemPhyXMini0246.acousticIntensityInWattsPerSquareMeter}
  \uses{def:physics:phyx-mini-0246:phyxminiproblems-problemphyxmini0246-acousticintensityquantity, def:physics:phyx-mini-0246:phyxminiproblems-problemphyxmini0246-acousticintensityreadout}
  SI watt-per-square-metre readout of an acoustic intensity
\end{definition}
\begin{proof}
  This is the declared model definition of acoustic Intensity In Watts Per Square Meter.
\end{proof}
\begin{definition}[Speaker Label]
  \label{def:physics:phyx-mini-0246:phyxminiproblems-problemphyxmini0246-speakerlabel}
  \lean{PhyXMiniProblems.ProblemPhyXMini0246.SpeakerLabel}
  The two loudspeakers shown on the left of the supplied figure
\end{definition}
\begin{proof}
  This is the declared model definition of Speaker Label.
\end{proof}
\begin{definition}[Figure Point]
  \label{def:physics:phyx-mini-0246:phyxminiproblems-problemphyxmini0246-figurepoint}
  \lean{PhyXMiniProblems.ProblemPhyXMini0246.FigurePoint}
  The four physical points distinguished in the planar model
\end{definition}
\begin{proof}
  This is the declared model definition of Figure Point.
\end{proof}
\begin{definition}[Ray Label]
  \label{def:physics:phyx-mini-0246:phyxminiproblems-problemphyxmini0246-raylabel}
  \lean{PhyXMiniProblems.ProblemPhyXMini0246.RayLabel}
  The ray labels printed on the source-to-B segments
\end{definition}
\begin{proof}
  This is the declared model definition of Ray Label.
\end{proof}
\begin{definition}[Acoustic Source Kind]
  \label{def:physics:phyx-mini-0246:phyxminiproblems-problemphyxmini0246-acousticsourcekind}
  \lean{PhyXMiniProblems.ProblemPhyXMini0246.AcousticSourceKind}
  The physical role of a source in the scenario
\end{definition}
\begin{proof}
  This is the declared model definition of Acoustic Source Kind.
\end{proof}
\begin{definition}[Planar Position]
  \label{def:physics:phyx-mini-0246:phyxminiproblems-problemphyxmini0246-planarposition}
  \lean{PhyXMiniProblems.ProblemPhyXMini0246.PlanarPosition}
  \uses{def:physics:phyx-mini-0246:phyxminiproblems-problemphyxmini0246-signedlengthquantity}
  A point in the source-observer plane with physical length coordinates
\end{definition}
\begin{proof}
  This is the declared model definition of Planar Position.
\end{proof}
\begin{definition}[Loudspeaker Interference Figure]
  \label{def:physics:phyx-mini-0246:phyxminiproblems-problemphyxmini0246-loudspeakerinterferencefigure}
  \lean{PhyXMiniProblems.ProblemPhyXMini0246.LoudspeakerInterferenceFigure}
  \uses{def:physics:phyx-mini-0246:phyxminiproblems-problemphyxmini0246-speakerlabel, def:physics:phyx-mini-0246:phyxminiproblems-problemphyxmini0246-figurepoint, def:physics:phyx-mini-0246:phyxminiproblems-problemphyxmini0246-raylabel, def:physics:phyx-mini-0246:phyxminiproblems-problemphyxmini0246-planarposition}
  Qualitative labels and placements read from the primary image. Metric readouts are supplied by MatchesSuppliedFigure; no ray length, phase at B, or total intensity is encoded in this structure
\end{definition}
\begin{proof}
  This is the declared model definition of Loudspeaker Interference Figure.
\end{proof}
\begin{definition}[Two Loudspeaker Setup]
  \label{def:physics:phyx-mini-0246:phyxminiproblems-problemphyxmini0246-twoloudspeakersetup}
  \lean{PhyXMiniProblems.ProblemPhyXMini0246.TwoLoudspeakerSetup}
  \uses{def:physics:phyx-mini-0246:phyxminiproblems-problemphyxmini0246-lengthquantity, def:physics:phyx-mini-0246:phyxminiproblems-problemphyxmini0246-acousticintensityquantity, def:physics:phyx-mini-0246:phyxminiproblems-problemphyxmini0246-speakerlabel, def:physics:phyx-mini-0246:phyxminiproblems-problemphyxmini0246-acousticsourcekind, def:physics:phyx-mini-0246:phyxminiproblems-problemphyxmini0246-loudspeakerinterferencefigure}
  The independent physical quantities and observables. AcousticAmplitude stays abstract because the problem asserts equal amplitudes without giving a unit or amplitude scale. The total intensity is an observable constrained only by the general interference law below
\end{definition}
\begin{proof}
  This is the declared model definition of Two Loudspeaker Setup.
\end{proof}
\begin{definition}[Matches Source Description]
  \label{def:physics:phyx-mini-0246:phyxminiproblems-problemphyxmini0246-matchessourcedescription}
  \lean{PhyXMiniProblems.ProblemPhyXMini0246.MatchesSourceDescription}
  \uses{def:physics:phyx-mini-0246:phyxminiproblems-problemphyxmini0246-twoloudspeakersetup}
  Source data from the prose and the Δφ₀ = 0 rad label: both objects are loudspeakers, have the common wavelength and equal amplitudes, start in phase, and each acting alone contributes the physical intensity denoted by I₀ at B. This contains no combined-intensity conclusion
\end{definition}
\begin{proof}
  This is the declared model definition of Matches Source Description.
\end{proof}
\begin{definition}[Matches Problem Readouts]
  \label{def:physics:phyx-mini-0246:phyxminiproblems-problemphyxmini0246-matchesproblemreadouts}
  \lean{PhyXMiniProblems.ProblemPhyXMini0246.MatchesProblemReadouts}
  \uses{def:physics:phyx-mini-0246:phyxminiproblems-problemphyxmini0246-lengthinmeters, def:physics:phyx-mini-0246:phyxminiproblems-problemphyxmini0246-twoloudspeakersetup}
  The wavelength, separation, and forward-distance readouts in metres
\end{definition}
\begin{proof}
  This is the declared model definition of Matches Problem Readouts.
\end{proof}
\begin{definition}[Matches Supplied Figure]
  \label{def:physics:phyx-mini-0246:phyxminiproblems-problemphyxmini0246-matchessuppliedfigure}
  \lean{PhyXMiniProblems.ProblemPhyXMini0246.MatchesSuppliedFigure}
  \uses{def:physics:phyx-mini-0246:phyxminiproblems-problemphyxmini0246-signedlengthinmeters, def:physics:phyx-mini-0246:phyxminiproblems-problemphyxmini0246-figurepoint, def:physics:phyx-mini-0246:phyxminiproblems-problemphyxmini0246-twoloudspeakersetup}
  Primary-image geometry in metre coordinates. The source plane is forward = 0; the speakers are at transverse coordinates ±3; A is at (10, 0) and B is at (10, 3), directly in front of the upper source. The horizontal and diagonal source-to-B rays are respectively labeled r₁ and r₂
\end{definition}
\begin{proof}
  This is the declared model definition of Matches Supplied Figure.
\end{proof}
\begin{definition}[Has Physical Acoustic Parameters]
  \label{def:physics:phyx-mini-0246:phyxminiproblems-problemphyxmini0246-hasphysicalacousticparameters}
  \lean{PhyXMiniProblems.ProblemPhyXMini0246.HasPhysicalAcousticParameters}
  \uses{def:physics:phyx-mini-0246:phyxminiproblems-problemphyxmini0246-lengthreadout, def:physics:phyx-mini-0246:phyxminiproblems-problemphyxmini0246-acousticintensityreadout, def:physics:phyx-mini-0246:phyxminiproblems-problemphyxmini0246-twoloudspeakersetup}
  Positivity and nondegeneracy of the physical quantities in the model
\end{definition}
\begin{proof}
  This is the declared model definition of Has Physical Acoustic Parameters.
\end{proof}
\begin{definition}[Satisfies Straight Line Ray Geometry]
  \label{def:physics:phyx-mini-0246:phyxminiproblems-problemphyxmini0246-satisfiesstraightlineraygeometry}
  \lean{PhyXMiniProblems.ProblemPhyXMini0246.SatisfiesStraightLineRayGeometry}
  \uses{def:physics:phyx-mini-0246:phyxminiproblems-problemphyxmini0246-lengthreadout, def:physics:phyx-mini-0246:phyxminiproblems-problemphyxmini0246-twoloudspeakersetup}
  Straight-line geometry for the two drawn rays. The upper source is aligned with B; the lower ray is the hypotenuse with legs equal to the forward distance and source separation. These are general relations, not numerical answers
\end{definition}
\begin{proof}
  This is the declared model definition of Satisfies Straight Line Ray Geometry.
\end{proof}
\begin{definition}[Satisfies Path Difference Law]
  \label{def:physics:phyx-mini-0246:phyxminiproblems-problemphyxmini0246-satisfiespathdifferencelaw}
  \lean{PhyXMiniProblems.ProblemPhyXMini0246.SatisfiesPathDifferenceLaw}
  \uses{def:physics:phyx-mini-0246:phyxminiproblems-problemphyxmini0246-lengthreadout, def:physics:phyx-mini-0246:phyxminiproblems-problemphyxmini0246-twoloudspeakersetup}
  The path difference is the magnitude of the two ray-length difference
\end{definition}
\begin{proof}
  This is the declared model definition of Satisfies Path Difference Law.
\end{proof}
\begin{definition}[Satisfies Propagation Phase Law]
  \label{def:physics:phyx-mini-0246:phyxminiproblems-problemphyxmini0246-satisfiespropagationphaselaw}
  \lean{PhyXMiniProblems.ProblemPhyXMini0246.SatisfiesPropagationPhaseLaw}
  \uses{def:physics:phyx-mini-0246:phyxminiproblems-problemphyxmini0246-lengthinmeters, def:physics:phyx-mini-0246:phyxminiproblems-problemphyxmini0246-twoloudspeakersetup}
  Propagation converts path difference to phase by 2π Δr / λ, in addition to the sources' initial phase difference. Radians are dimensionless
\end{definition}
\begin{proof}
  This is the declared model definition of Satisfies Propagation Phase Law.
\end{proof}
\begin{definition}[Satisfies Coherent Two Source Interference Law]
  \label{def:physics:phyx-mini-0246:phyxminiproblems-problemphyxmini0246-satisfiescoherenttwosourceinterferencelaw}
  \lean{PhyXMiniProblems.ProblemPhyXMini0246.SatisfiesCoherentTwoSourceInterferenceLaw}
  \uses{def:physics:phyx-mini-0246:phyxminiproblems-problemphyxmini0246-acousticintensityreadout, def:physics:phyx-mini-0246:phyxminiproblems-problemphyxmini0246-twoloudspeakersetup}
  The general coherent two-source intensity law. In any coherent unit system, contributions I₁ and I₂ with received phase difference Δφ combine as I₁ + I₂ + 2 sqrt(I₁ I₂) cos(Δφ). No problem-specific path or answer coefficient occurs in this law
\end{definition}
\begin{proof}
  This is the declared model definition of Satisfies Coherent Two Source Interference Law.
\end{proof}
\begin{definition}[intensity Ratio At B]
  \label{def:physics:phyx-mini-0246:phyxminiproblems-problemphyxmini0246-intensityratioatb}
  \lean{PhyXMiniProblems.ProblemPhyXMini0246.intensityRatioAtB}
  \uses{def:physics:phyx-mini-0246:phyxminiproblems-problemphyxmini0246-acousticintensityinwattspersquaremeter, def:physics:phyx-mini-0246:phyxminiproblems-problemphyxmini0246-twoloudspeakersetup}
  The total intensity at B as a dimensionless multiple of I₀
\end{definition}
\begin{proof}
  This is the declared model definition of intensity Ratio At B.
\end{proof}
\begin{definition}[rounded To Nearest Hundredth]
  \label{def:physics:phyx-mini-0246:phyxminiproblems-problemphyxmini0246-roundedtonearesthundredth}
  \lean{PhyXMiniProblems.ProblemPhyXMini0246.roundedToNearestHundredth}
  Round a dimensionless scalar to the nearest hundredth
\end{definition}
\begin{proof}
  This is the declared model definition of rounded To Nearest Hundredth.
\end{proof}
\begin{definition}[Answer Choice]
  \label{def:physics:phyx-mini-0246:phyxminiproblems-problemphyxmini0246-answerchoice}
  \lean{PhyXMiniProblems.ProblemPhyXMini0246.AnswerChoice}
  The answer labels printed with the problem
\end{definition}
\begin{proof}
  This is the declared model definition of Answer Choice.
\end{proof}
\begin{definition}[displayed Intensity Multiple]
  \label{def:physics:phyx-mini-0246:phyxminiproblems-problemphyxmini0246-displayedintensitymultiple}
  \lean{PhyXMiniProblems.ProblemPhyXMini0246.displayedIntensityMultiple}
  \uses{def:physics:phyx-mini-0246:phyxminiproblems-problemphyxmini0246-answerchoice}
  The displayed coefficient multiplying I₀ for each answer choice
\end{definition}
\begin{proof}
  This is the declared model definition of displayed Intensity Multiple.
\end{proof}
\begin{definition}[Matches Displayed Intensity Choice]
  \label{def:physics:phyx-mini-0246:phyxminiproblems-problemphyxmini0246-matchesdisplayedintensitychoice}
  \lean{PhyXMiniProblems.ProblemPhyXMini0246.MatchesDisplayedIntensityChoice}
  \uses{def:physics:phyx-mini-0246:phyxminiproblems-problemphyxmini0246-twoloudspeakersetup, def:physics:phyx-mini-0246:phyxminiproblems-problemphyxmini0246-intensityratioatb, def:physics:phyx-mini-0246:phyxminiproblems-problemphyxmini0246-roundedtonearesthundredth, def:physics:phyx-mini-0246:phyxminiproblems-problemphyxmini0246-answerchoice, def:physics:phyx-mini-0246:phyxminiproblems-problemphyxmini0246-displayedintensitymultiple}
  A choice agrees with the derived intensity ratio to the nearest hundredth
\end{definition}
\begin{proof}
  This is the declared model definition of Matches Displayed Intensity Choice.
\end{proof}
\begin{definition}[Is Unique Matching Displayed Intensity Choice]
  \label{def:physics:phyx-mini-0246:phyxminiproblems-problemphyxmini0246-isuniquematchingdisplayedintensitychoice}
  \lean{PhyXMiniProblems.ProblemPhyXMini0246.IsUniqueMatchingDisplayedIntensityChoice}
  \uses{def:physics:phyx-mini-0246:phyxminiproblems-problemphyxmini0246-twoloudspeakersetup, def:physics:phyx-mini-0246:phyxminiproblems-problemphyxmini0246-answerchoice, def:physics:phyx-mini-0246:phyxminiproblems-problemphyxmini0246-matchesdisplayedintensitychoice}
  Exactly one displayed coefficient matches the rounded derived ratio
\end{definition}
\begin{proof}
  This is the declared model definition of Is Unique Matching Displayed Intensity Choice.
\end{proof}
\begin{lemma}[ray lengths path difference and phase at B]
  \label{lem:physics:phyx-mini-0246:phyxminiproblems-problemphyxmini0246-ray-lengths-path-difference-and-phase-at-b}
  \lean{PhyXMiniProblems.ProblemPhyXMini0246.ray_lengths_path_difference_and_phase_at_B}
  \uses{def:physics:phyx-mini-0246:phyxminiproblems-problemphyxmini0246-lengthinmeters, def:physics:phyx-mini-0246:phyxminiproblems-problemphyxmini0246-twoloudspeakersetup, def:physics:phyx-mini-0246:phyxminiproblems-problemphyxmini0246-matchessourcedescription, def:physics:phyx-mini-0246:phyxminiproblems-problemphyxmini0246-matchesproblemreadouts, def:physics:phyx-mini-0246:phyxminiproblems-problemphyxmini0246-hasphysicalacousticparameters, def:physics:phyx-mini-0246:phyxminiproblems-problemphyxmini0246-satisfiesstraightlineraygeometry, def:physics:phyx-mini-0246:phyxminiproblems-problemphyxmini0246-satisfiespathdifferencelaw, def:physics:phyx-mini-0246:phyxminiproblems-problemphyxmini0246-satisfiespropagationphaselaw}
  The ray geometry gives r₁ = 10 m, r₂ = sqrt 136 m, path difference sqrt 136 - 10 m, and the displayed received phase. All four statements are derived conclusions rather than setup assumptions
\end{lemma}
\begin{proof}
  The conclusion follows from the governing relations and figure readouts named in the statement of ray lengths path difference and phase at B.
\end{proof}
\begin{lemma}[exact intensity at B]
  \label{lem:physics:phyx-mini-0246:phyxminiproblems-problemphyxmini0246-exact-intensity-at-b}
  \lean{PhyXMiniProblems.ProblemPhyXMini0246.exact_intensity_at_B}
  \uses{def:physics:phyx-mini-0246:phyxminiproblems-problemphyxmini0246-acousticintensityreadout, def:physics:phyx-mini-0246:phyxminiproblems-problemphyxmini0246-twoloudspeakersetup, def:physics:phyx-mini-0246:phyxminiproblems-problemphyxmini0246-matchessourcedescription, def:physics:phyx-mini-0246:phyxminiproblems-problemphyxmini0246-matchesproblemreadouts, def:physics:phyx-mini-0246:phyxminiproblems-problemphyxmini0246-hasphysicalacousticparameters, def:physics:phyx-mini-0246:phyxminiproblems-problemphyxmini0246-satisfiesstraightlineraygeometry, def:physics:phyx-mini-0246:phyxminiproblems-problemphyxmini0246-satisfiespathdifferencelaw, def:physics:phyx-mini-0246:phyxminiproblems-problemphyxmini0246-satisfiespropagationphaselaw, def:physics:phyx-mini-0246:phyxminiproblems-problemphyxmini0246-satisfiescoherenttwosourceinterferencelaw}
  Substitution of the equal single-source intensities into the coherent-source law yields the exact, unrounded intensity factor in every coherent unit system
\end{lemma}
\begin{proof}
  The conclusion follows from the governing relations and figure readouts named in the statement of exact intensity at B.
\end{proof}
\begin{lemma}[intensity ratio at B exact]
  \label{lem:physics:phyx-mini-0246:phyxminiproblems-problemphyxmini0246-intensity-ratio-at-b-exact}
  \lean{PhyXMiniProblems.ProblemPhyXMini0246.intensity_ratio_at_B_exact}
  \uses{def:physics:phyx-mini-0246:phyxminiproblems-problemphyxmini0246-twoloudspeakersetup, def:physics:phyx-mini-0246:phyxminiproblems-problemphyxmini0246-matchessourcedescription, def:physics:phyx-mini-0246:phyxminiproblems-problemphyxmini0246-matchesproblemreadouts, def:physics:phyx-mini-0246:phyxminiproblems-problemphyxmini0246-hasphysicalacousticparameters, def:physics:phyx-mini-0246:phyxminiproblems-problemphyxmini0246-satisfiesstraightlineraygeometry, def:physics:phyx-mini-0246:phyxminiproblems-problemphyxmini0246-satisfiespathdifferencelaw, def:physics:phyx-mini-0246:phyxminiproblems-problemphyxmini0246-satisfiespropagationphaselaw, def:physics:phyx-mini-0246:phyxminiproblems-problemphyxmini0246-satisfiescoherenttwosourceinterferencelaw, def:physics:phyx-mini-0246:phyxminiproblems-problemphyxmini0246-intensityratioatb}
  The exact dimensionless intensity ratio at point B
\end{lemma}
\begin{proof}
  The conclusion follows from the governing relations and figure readouts named in the statement of intensity ratio at B exact.
\end{proof}
% --- Archon declaration topology end ---
% --- Archon physics formalization source end ---
